% Diagramme de classes de la hiérarchie des agents (figure 3.5)
\begin{tikzpicture}[x=1cm, y=1cm]

  \tikzset{
    cl/.style={draw=black, fill=white, rectangle split, rectangle split parts=3,
               rectangle split draw splits=true, text width=52mm,
               inner sep=2.6pt, font=\sffamily\tiny},
    cl2/.style={draw=black, fill=white, rectangle split, rectangle split parts=2,
                rectangle split draw splits=true, text width=36mm,
                inner sep=2.6pt, font=\sffamily\tiny},
    pers/.style={fbox, minimum width=26mm, minimum height=9mm,
                 font=\sffamily\tiny, align=center}
  }
  \newcommand{\nomcl}[1]{\makebox[\linewidth][c]{\textbf{#1}}}
  \newcommand{\stereo}[1]{\makebox[\linewidth][c]{$\ll$ #1 $\gg$}}

  % ── Classe de base ────────────────────────────────────────
  \node[cl] (base) at (0,3.4) {
    \stereo{abstraite}

    \nomcl{AgentConversationnelDeBase}
    \nodepart{two}
      \textminus\ instructions : Texte\\
      \textminus\ capacites : Capacite [0..*]\\
      \textminus\ languePreferee : Langue
    \nodepart{three}
      + surTourClientTermine(tour)\\
      + evaluerSentiment(tour) : Sentiment\\
      + journaliserTour(tour)\\
      + cloturerConversation()\\
      + changerLangue(langue)\\
      \# regleAbstention() : Texte
  };

  % ── Capacités détenues ────────────────────────────────────
  \node[cl2, text width=33mm] (cap) at (6.7,3.4) {
    \nomcl{Capacite}
    \nodepart{two}
      \textminus\ nom : Texte\\
      \textminus\ estSensible : Booleen\\
      \textminus\ estDeTransfert : Booleen
  };

  \draw[agregation] (base.east) -- node[card, above, pos=0.16] {1}
                    node[card, above, pos=0.86] {0..*} (cap.west);

  % ── Personas ──────────────────────────────────────────────
  \node[pers] (p1) at (-5.9,-1.6) {AgentAccueil};
  \node[pers] (p2) at (-2.95,-1.6) {AgentOrientation};
  \node[pers] (p3) at ( 0.0,-1.6) {AgentFacturation};
  \node[pers] (p4) at ( 2.95,-1.6) {AgentTechnique};
  \node[pers] (p5) at ( 5.9,-1.6) {AgentGestion\\DeCompte};

  % Généralisation en arbre
  \draw[heritage] (0,0.50) -- (base.south);
  \draw (-5.9,0.50) -- (5.9,0.50);
  \foreach \p in {p1,p2,p3,p4,p5}{ \draw (\p.north) -- (\p.north |- 0,0.50); }

  % ── Domaine de spécialité ─────────────────────────────────
  \node[cl2] (dom) at (2.95,-5.1) {
    \nomcl{DomaineDeSpecialite}
    \nodepart{two}
      \textminus\ cle : Texte\\
      \textminus\ sujetsCouverts : Texte\\
      \textminus\ capaciteDeTransfert : Texte\\
      \textminus\ phraseDeTransition : Map[Langue]
  };

  % Seuls les trois agents spécialisés couvrent un domaine.
  \foreach \p in {p3,p4,p5}{
    \draw[dirige] (\p.south) -- node[card, right, pos=0.75] {1} (dom.north);
  }

  % ── Contraintes structurantes ─────────────────────────────
  \node[note, text width=50mm, anchor=north east] at (-2.05,-3.35)
       {\{ contrainte \} Les instructions d'un agent sont dérivées de l'ensemble de ses capacités. Un agent ne peut donc pas recevoir la consigne d'utiliser une capacité qu'il ne possède pas.};
  \draw[dependance] (-4.55,-3.35) -- (-4.35,-2.05);

  \node[note, text width=42mm, anchor=north west] at (5.55,-3.55)
       {La phrase de transition est fixée pour chaque langue, et non produite par le modèle.};
  \draw[dependance] (5.55,-4.35) -- (4.75,-4.85);

\end{tikzpicture}
